% RQ5 Discussion: Inter-DAO Cooperation

\subsection{Interpretation of Results}

\subsubsection{The Complementarity Paradox}

Our results reveal a fundamental tension in inter-DAO cooperation: the conditions that make cooperation most valuable (heterogeneous, complementary organizations) also make it hardest to achieve.

Homogeneous DAOs share preferences and norms, making agreement easy. But they lack complementary capabilities, so cooperation adds limited value---they're essentially duplicating each other.

Heterogeneous DAOs bring diverse resources and perspectives, creating potential for substantial value creation. But preference divergence makes agreement difficult, and size differences create bargaining asymmetries.

This suggests that DAOs seeking high-value cooperation should:
\begin{itemize}
    \item Invest in relationship building before proposing joint initiatives
    \item Design proposals that create clear value for all parties
    \item Use explicit fairness mechanisms to address asymmetries
    \item Start with smaller, trust-building collaborations before major initiatives
\end{itemize}

\subsubsection{Fairness as Infrastructure}

Our asymmetric scenario results underscore that fairness mechanisms are not optional---they are infrastructure for sustainable cooperation.

Without explicit fairness guarantees, small DAOs rationally distrust partnerships with larger organizations. The risk of domination or exploitation outweighs potential benefits.

Shapley-style division, which allocates benefits proportional to marginal contribution, provides theoretical justification for fair sharing. But implementation requires:
\begin{itemize}
    \item Clear value measurement
    \item Transparent accounting of contributions
    \item Dispute resolution mechanisms
    \item Exit rights for dissatisfied parties
\end{itemize}

\subsubsection{The Hub Advantage}

Hub-spoke structures significantly improve multi-party coordination. A central coordinator reduces bilateral negotiation costs from $O(n^2)$ to $O(n)$, enabling larger coalitions.

However, hub structures create their own risks:
\begin{itemize}
    \item Hub capture: Central coordinator accumulates disproportionate influence
    \item Single point of failure: Coalition depends on hub functioning
    \item Spoke dependency: Spoke DAOs may lose autonomy
\end{itemize}

Mitigation strategies include:
\begin{itemize}
    \item Hub rotation or election
    \item Spoke veto rights on hub actions
    \item Transparency requirements for hub operations
\end{itemize}

\subsection{Design Principles for Inter-DAO Cooperation}

Based on our results, we propose the following principles:

\begin{enumerate}
    \item \textbf{Start small}: Begin with low-stakes collaborations to build trust
    \item \textbf{Explicit fairness}: Use clear, defensible division rules (Shapley, proportional)
    \item \textbf{Overlapping membership}: Encourage cross-DAO participation to align incentives
    \item \textbf{Clear value proposition}: Each party should understand their specific benefit
    \item \textbf{Exit rights}: Maintain ability to withdraw from cooperation
    \item \textbf{Coordination infrastructure}: For multi-party cooperation, designate coordinator role
\end{enumerate}

\subsection{Comparison with Real Inter-DAO Cooperation}

\subsubsection{Protocol Guild}

Protocol Guild's multi-DAO funding of Ethereum developers exemplifies successful heterogeneous cooperation. Key features matching our findings:
\begin{itemize}
    \item Clear shared interest (Ethereum infrastructure)
    \item Transparent contribution and distribution
    \item Low coordination overhead (passive contribution)
\end{itemize}

\subsubsection{Metagovernance}

DAOs holding and voting tokens of other DAOs creates structural overlap. Our simulations support this approach as facilitating cooperation.

\subsubsection{Gitcoin Grants}

Quadratic funding with matching from multiple treasuries demonstrates coordinated public goods provision. The mechanism design addresses free-rider problems our simulations identify.

\subsection{Limitations}

\subsubsection{Simplified Preference Model}

We model DAO preferences as vectors, but real preferences are:
\begin{itemize}
    \item Multi-dimensional and complex
    \item Evolving over time
    \item Subject to internal politics
\end{itemize}

\subsubsection{No Repeated Game Dynamics}

Our simulations run for fixed periods. Long-term repeated interactions may:
\begin{itemize}
    \item Build trust beyond what we capture
    \item Enable reputation-based cooperation
    \item Create path dependencies
\end{itemize}

\subsubsection{Exogenous Cooperation Opportunities}

We generate cooperation opportunities randomly. In practice, opportunities are endogenous---they emerge from relationships and shared challenges.

\subsection{Future Work}

\begin{enumerate}
    \item \textbf{Repeated game dynamics}: Model long-term inter-DAO relationships with reputation effects
    \item \textbf{Network formation}: Endogenize cooperation structure evolution
    \item \textbf{Adversarial scenarios}: Model attempts to exploit inter-DAO mechanisms
    \item \textbf{Empirical validation}: Compare simulation predictions to observed inter-DAO cooperation
\end{enumerate}
